% ===========================================================================
% 05-Limitations.tex. PRESENTATION ONLY. Every number, interval, hedge and
% conclusion is verbatim from v1. ONE \needsaudit, in "Cross-plate absolutes".
% ===========================================================================

\chapter{Limitations and future research directions}
\label{sec:limitations}
\framedecl{ER}

\begin{question}
Every measurement in this thesis was made on one atlas, one checkpoint and one
slice of both. Which of its conclusions are properties of the model, and which
are properties of that slice?
\end{question}

\noindent
A limitation is worth writing down when it changes what a number licenses.

% ===========================================================================
\section{Every result here is scoped, and the scopes differ}

\paragraph{Drug coverage} A study of $106$ drugs can establish that a model uses
the drug it is handed; it cannot describe how drugs in general behave. The
residual analysis uses $\num{620564}$ cells covering $106$ drugs across $50$
cell lines, roughly a tenth of the atlas's drug set, ample for testing
whether the model uses the drug and for measuring transfer but not a population
estimate. Claims about the \emph{distribution} of drug behaviours belong to the
broader characterisation, which used $\num{6628}$ conditions.

\paragraph{Only reproducible conditions} A residual that does not reproduce
between two halves of the same well is not a target anything could learn, so the
residual arm trains on the $50.5\%$ of conditions whose residual does reproduce,
$\num{2523}$ of $\num{4999}$ training conditions kept at $\cos > 0.109$, the
$95$th percentile of a null pairing half A of one condition with half B of
another.\gloss{reliability line}{The half-split cosine threshold above which a
condition's residual counts as reproducible.} The model was taught the
empirically reproducible fraction. Split-half retrieval establishes
discriminability of treated populations under this sampling design and cell
budget; it does not establish predictability from untreated input or
learnability.

A reliability filter might be expected to select \emph{toward}
context-in\-de\-pen\-dent conditions and bias the transfer coefficient upward.
It does the opposite. Unfiltered \ci{0.597}{0.563}{0.632}, filtered
\ci{0.553}{0.514}{0.593}; the filtered value is quoted, because disattenuation
is unstable at the very low reliabilities the filter removes.

That calculation's own filter is the inherited $\cos > 0.2$ line, where the
training set uses $0.109$, so the two estimates bracket the trained-on
population without reproducing it. The generation evaluation mismatches
differently again, applying $0.109$ to \texttt{train} and leaving the held-out
splits unfiltered. That one is load-bearing and is taken up below.

\paragraph{The intervals omit generation variance} The intervals here are
cluster-robust with each condition's score treated as fixed: two-way over
\emph{cell line} and \emph{treated well} on a $\min(G_{\textrm{line}},
G_{\textrm{well}}) - 1$ Student $t$ reference for the headline gaps, reported
again clustered on the \emph{drug}; multiway dyadic for the transfer
coefficient; a one-way clustered bootstrap, unit named, in a few tables.

They capture between-cluster variance and leave out the variance from sampling
the model's predictions. With four sampled generations per condition at temperature
$0.8$ in the residual-frame arms and eight in the expression-frame ones
(\cref{sec:methods-data}), that source is not negligible, and no across-seed
replication of the generation arms was run, deferred or planned. Every interval
quoted from a generation arm covers the sampling of conditions, cell lines and
wells, and stops there.

One measurement \emph{was} repeated across seeds, recorded so it is not mistaken
for the replication that was not run. The activation probe of
\cref{sec:res-mechanism} returned $+0.0197$, $+0.0118$ and $+0.0085$ on three
seeds, each excluding zero. Its spread does not transfer to $\NIR$ and does not
establish a scale for generation variance. Those runs share a checkpoint and
score teacher-forced log-probabilities, so their seed moves the prompts, the
partners and the bootstrap, and leaves both the weights and the sampled
prediction untouched. They are scored on $60$, $240$ and $240$ prompts in that order, so
the largest of the three effects is also the smallest-$n$ one.

\paragraph{No training run was repeated} The seeds above move prompts,
partners and resampling. None of them moves the weights. Every arm here was
trained once, from one base checkpoint under identical hyperparameters
(\cref{tab:trainconfig}), so each dissociation this thesis rests on --- the
residual re-encoding against the standard target, and the optimal-transport
target against both --- compares one training run with one training run. A
second run of the same arm under a different data order would land on a
different checkpoint and score differently, and the size of that difference was
never measured here. It is therefore not covered by any interval in this
document: the intervals quantify sampling over conditions, cell lines and wells,
and a between-arm gap of the size reported would survive or fail a training
replication that nobody has run. The ten-epoch diagnostic varies the horizon and
not the seed, so it bounds undertraining and not run-to-run spread. What this
costs is specific rather than general: the \emph{direction} of the re-encoding
effect is corroborated by the token-divergence measurement, which needs no
model at all, but the \emph{magnitude} of every arm-to-arm difference is an
n-of-one comparison.

\paragraph{A swing that looked like generation variance and was not} Unseeded
generation is the kind of defect that surfaces as an unstable number, and once
it seemed to.

\begin{withdrawn}{seed instability in the held-out-drug control}
Two superseded evaluations of the held-out-drug control returned values of
$\gap$ far enough apart to be recorded as evidence of seed instability. Only one
of the pair survives in the results archive, so the size of the swing is no
longer quotable. The diagnosis was wrong. It named the trigger and missed the
cause.
\end{withdrawn}

Both of the superseded figures are gaps against the \texttt{opposite} swap, which
\cref{sec:res-generalisation} shows is not a null; it carries a term for how far
the \emph{partner} drug was pushed the wrong way, which changes with the
draw.\seealso{\cref{sec:res-generalisation}, where a neutral stratum replaces
\texttt{opposite} as the comparator} On the canonical evaluation the same
control reads $+0.0379$ against \texttt{opposite} and
\ci{-0.0315}{-0.0836}{+0.0206} against the neutral stratum, clustered two-way on
cell line and well.

Two mechanisms were at work, and both are now addressed. \texttt{do\_sample}
advanced from global torch state, so the two arms of a contrast were drawn from
different points of one
stream; each call now seeds its own generator from (condition, arm, batch), and
the comparator has been corrected. The general caution stands regardless; this
swing is no longer evidence for it.

\paragraph{Pseudobulk resolution} Residuals are condition-level pseudobulks over
$\geq 40$ treated cells, so single-cell heterogeneity of the response is not
modelled in this frame. Earlier project results show why: at single-cell
resolution the drug moves a cell less than noise does (SNR $\approx 0.75$,
median $0$ differentially expressed genes at FDR $q<0.05$).

\paragraph{Cross-plate absolutes} Residual-space scores use comparison sets
grouped by cell line, not by plate. The \emph{difference} $\gap$ is leak-immune,
since both arms share the same control cell. The absolute $\NIR$ values are a
different matter.

\cref{sec:methods-residual} argues the grouping is admissible because the
plate-scoped generic has already subtracted each plate's mean response from
truths and predictions alike. But it removes plate structure in expectation
only, and the generic is estimated from a small plate group, so some may survive
in an individual residual. The rescore that would size what survives, regrouping
the generation-free arms within plate, needs no GPU time and was not run.

What exists instead is one-sided. \cref{tab:plateleak} bounds it from
above.\framecue{\Eframe}{the plate-leak audit} In the expression frame with no
generic subtracted, holding the plate constant moves a drug-agnostic
\needsaudit[same construction as the DRF negative control]{mean baseline} from
$0.399$ to $0.161$ and the precision reference from $0.625$ to $0.575$ --- the
yardstick falls with the baseline (\cref{fig:plateleak}), so the leak is not an
offset one side of a comparison could be corrected by.
Residual-frame absolutes should be read within their own frame and never set
beside a within-plate number.

\begin{table}[htbp]
\centering
\begin{threeparttable}
\caption{\textbf{Expression-frame $\NIR$, the leak audit.} A drug-agnostic mean
baseline more than doubles its score when the plate is free to vary. The
precision reference falls once the plate is held constant, because batch
structure was inflating it too.}
\label{tab:plateleak}
\footnotesize
\setlength{\tabcolsep}{6pt}
\begin{tabular}{@{}lrr@{}}
\toprule
\thead{Predictor (drug-agnostic)} & \thead{cross-plate}
  & \thead{within plate} \\
\midrule
mean baseline & 0.399 & \textbf{0.161} \\
control-copy & 0.551 & 0.510 \\
\midrule
precision reference (within-well split-half) & 0.625 & 0.575 \\
fraction of drugs at chance & 43\% & 49\% \\
\bottomrule
\end{tabular}
\begin{tablenotes}[flushleft]\footnotesize
\item A wider diagnostic run than the $n = 606$ tier-2 benchmark the blindness
verdict is stated on, at $n = \num{1163}$ cross-plate and $n = 998$ within
plate. Holding the plate constant costs every condition with no same-plate
comparison set, so the within-plate column is a subset of the cross-plate one
rather than the same conditions rescored.
\item Its within-plate figures therefore sit slightly off \cref{tab:allarms},
control-copy $0.510$ against $0.504$, mean baseline $0.161$ against $0.180$,
precision reference $0.575$ against $0.576$.
\end{tablenotes}
\end{threeparttable}
\end{table}

\begin{figure}[htbp]
\centering
% drawn at 142mm by thesis_v2/figs/plateleak.py, the width of the measure; no
% width= key, so it is placed at scale 1.0 and its type prints at the point
% sizes it was set in. Values in figs/fig-plateleak.json, read from
% RESULTS_cluster/leak_crossplate.json and RESULTS_cluster/leak_sameplate.json,
% tiers.tier2_unseen_drugs.agg, the *_nir_expr entries only.
\includegraphics{fig-plateleak}
\caption[Holding the plate constant lowers the reference too]{\textbf{The batch
shortcut inflates the precision reference as well as the baseline, so it is not
an offset that could be subtracted from one side of a comparison.}
Expression-frame $\NIR$, tier 2 unseen drugs, drug-agnostic arms only: the same
audit as \cref{tab:plateleak}, with the linear map read from the same block of
the same two records. Forcing the comparison set within plate collapses a mean
baseline that carries no drug information whatsoever, $0.399$ to $0.161$, and
leaves control-copy close to where it was, $0.551$ to $0.510$ --- but it also
lowers the within-well split-half precision reference, $0.625$ to $0.575$, by
about as much as it lowers control-copy. That reference is the yardstick these
expression-frame numbers are read against, so the leak has no one side to be
corrected on. The two columns are not the same conditions: holding the plate
constant costs every condition with no same-plate comparison set, so the
within-plate column is a \emph{subset} --- $n = 998$ conditions, $46$ drugs,
$41$ cell lines --- of the cross-plate column's $n = \num{1163}$, $50$ drugs and
$43$ cell lines, and each segment joins two scorings of overlapping but unequal
populations rather than one set of conditions rescored. The marks are bare
because neither record carries an interval for these quantities, each entry
being a value and its $n$, and this document draws only $95\%$ bootstrap
intervals clustered over cell lines. The rank-frame variants stored beside these
values are a different frame and are not drawn, and none of the values here may
be set beside a residual-frame number.}
\label{fig:plateleak}
\end{figure}

% "Frame coverage". The shares are SQUARED EUCLIDEAN NORMS, not variances:
% nothing is centred anywhere in the computation, so ||s||^2 is a sum of squares
% about the origin. shared/inference.py::energy_shares calls it an "exact
% variance decomposition" in its docstring -- that docstring is wrong and the
% prose must not follow it. Construction: shift = residual + generic, so
%     ||s||^2 = ||r||^2 + ||g||^2 + 2<r,g>
% and the three shares sum to one by construction, cross term included.
% Values from RESULTS_cluster/vardecomp_matched.json, main.scope:
% residual_energy_share 0.62112, generic_energy_share 0.38993,
% cross_energy_share -0.01106, energy_shares_sum 1.000. The 0.786 and 0.620 are
% the separate norm RATIOS residual_over_shift and generic_over_shift, NOT the
% square roots of the shares: variance_decomposition.py:632-634 takes a mean over
% conditions in both cases, so 0.786^2 = 0.617, not 0.621. The final sentence
% exists to stop a reader squaring and finding a discrepancy, and it must name
% BOTH sides: the shares fail to be the squares of the ratios only because both
% quantities are averaged over conditions and squaring does not commute with
% averaging. Asserting it of the shares alone does not entail the conclusion.
% For the same reason 0.786 is set as prose and not as ||r||/||s|| = 0.786: the
% norm entry in the Symbols list of 00-HowToRead.tex fixes ||.|| as the
% Euclidean norm of a vector, so that equation would claim a ratio of two norms,
% while the estimator is a mean over conditions of per-condition ratios. The
% JSON keys keep the word "energy"; the prose does not, and that same norm entry
% is what defines the term used here.
\paragraph{Frame coverage} The drug-specific residual is on average $0.786$ of
the shift in norm, against $0.620$ for the generic programme and batch, and the
two are essentially orthogonal, the shift's squared norm decomposing into
residual $0.621$, generic $0.390$ and a cross term of $-0.011$, summing to
$1.000$. Shares and ratios are both averaged over conditions, so the shares are
not the squares of the two ratios above. The frame covers
$62.1\%$ of the shift's \emph{squared norm}, a real but not total share, and
every result here should be read within that bound. The generic $39.0\%$ is
drug-agnostic and trivially matched, part of what a drug-agnostic predictor
supplies on absolute-accuracy metrics, though the generic programme is not what
saturates $\DEdr$; reversion toward a central point is (\cref{sec:res-metrics}).

% "The similarity is a choice too". Pairs with "Frame coverage" above: that
% paragraph states what the frame's choice of scored OBJECT costs, this one what
% its choice of SIMILARITY costs. Sources, all read directly:
%   nir_benchmark.py:13     "expr-NIR : Euclidean distance after decoding ranks
%                            -> expression via linear_model.json"
%   nir_benchmark.py:363-4  d_own = np.linalg.norm(pe - truth_expr[d]) etc.
%   nir_benchmark.py:864    "expression_distance": "Euclidean"
%   residual_eval.py:13     "Similarity = cosine."
%   residual_eval.py:125    def cos(a, b)
% The rescore: RESULTS_cluster/nir_sameplate_profiles.npz, 606 predictions and
% truths over 80 groups, tier2_unseen_drugs, the same support as the frozen
% manifest, scored under the definition of sec:res-metrics with the half-tie
% rule. -||.||_2 gives mean 0.4977 median 0.5000; cosine gives mean 0.5007
% median 0.5000; Spearman 0.6217 between the two per condition; 352/606 move by
% more than 0.10. The Euclidean value reproduces the published 0.498 on the same
% n, which is the only thing validating the rescore, so it is stated.
% The Spearman was re-derived from the same npz under the same half-tie rule
% (scipy spearmanr on the 606 per-condition pairs; Pearson 0.6218, Kendall
% 0.4785 on the same pairs). Restricting to groups of at least three drugs and
% dropping the half-tie rule both return 0.6217; averaging Spearman within
% groups returns 0.5775, a different estimator and not the one quoted.
% NOT RESCORABLE: every local .npz is expression-frame. There is no
% re_*_profiles.npz, so the residual frame carries no such measurement and the
% last sentence says so, which is what stops a reader generalising.
% The sigma instances belong at the definition in 04a-ruler.tex, which already
% parameterises by sigma and already declares the expression instance; this pass
% was scoped to one file, so the residual instance is named here instead.
\paragraph{The similarity is a choice too} $\NIR$ is parameterised by a
similarity $\sigma$ (\cref{sec:res-metrics}), and the two frames instantiate it
differently: the expression frame scores $\sigma(\hat y, y) = -\lVert \hat y - y
\rVert$, the residual frame the cosine. What that choice carries was measured on
the run the blindness verdict is stated on. The same $606$ stored predictions
and truths, the same $80$ comparison groups, tier 2 within plate, were rescored
under both similarities using the definition above, half-tie rule included: the
aggregate reads $0.4977$ under the negated Euclidean distance against $0.5007$
under the cosine, both at chance. The Euclidean value reproduces the published
figure on that support, which is what validates the rescore.

What does not hold across the substitution is the ordering beneath the
aggregate. The two scorings agree per condition at Spearman $0.6217$, and $352$
of the $606$ conditions move by more than $0.10$. That neither strengthens nor
weakens the verdict. A predictor at chance leaves no signal for either
similarity to find, so what separates the two orderings is noise, and the
aggregate stands under both.

Two properties of the published convention belong beside it. A Euclidean
distance is dominated by its largest coordinates, so the high-abundance genes
carry the comparison. And the vectors compared are reconstructions: the
prediction is a cell sentence, decoded to expression through a fitted linear map
before any distance is taken, and that map was fitted for absolute fidelity and
never validated for discrimination (\cref{sec:lit-c2s}). The residual frame
carries no such measurement, because its predictions and truths were not stored
in a form that can be rescored, so nothing is claimed about its sensitivity to
$\sigma$ in either direction.

\begin{scopelimit}[carried into \cref{sec:res-blind} and every per-condition
reading of an expression-frame score]
An at-chance aggregate $\NIR$ in the expression frame holds under either
similarity, negated Euclidean or cosine, on the one run rescored here, tier 2
within plate. Expression-frame aggregates that carry signal were not rescored,
so their standing under $\sigma$ is untested. No per-condition or per-stratum
reading has that standing either: on the run that carries the verdict the two
orderings agree only at Spearman $0.6217$, so no figure below the aggregate may
be presented as independent of $\sigma$. The residual frame carries no such
measurement, so its scores are covered neither way.
\end{scopelimit}

\paragraph{The splits are scored on different populations} The within-well
split-half precision reference differs sharply across the three splits as
scored, $0.956$ on \texttt{train}, $0.824$ on \texttt{unseen\_combo} and $0.836$
on \texttt{unseen\_drug}. The obvious reading is that held-out conditions carry
less recoverable signal. That reading is wrong.

It is the selection asymmetry noted above. \texttt{train} is drawn from a pool
already filtered at the $0.109$ line, and none of its $300$ scored conditions
falls at or below it, while the held-out splits are unfiltered and still contain
conditions whose two halves disagree outright. Impose the training line on all
three and the references agree at $0.955$, $0.961$ and $0.960$, a spread of
$0.006$ where there had been $0.132$ (\cref{fig:splitref}). Matched on
reproducibility, a held-out condition is as recoverable as a trained one.

% The float that stood here (193-249 of the previous revision) is now
% figs/splitref.tex. It moved out of the section for the same reason the
% other reworked figures did -- the drawing carries 70 lines of its own
% rationale, which does not belong in the prose -- and it lost its
% \begin{fullwidth} on the way: it reserved 174mm, drew 150.2mm, and cost
% this page its margin column to show six numbers the caption already states.
% It now draws 132.9mm inside the 142mm measure.
% ===========================================================================
% splitref.tex -- the within-well split-half reference, as the three splits are
% actually scored and with the resolved +0.1086 reliability line imposed on all
% three.
% \input by Sections/05-Limitations.tex. Compiles against thesis_v2/preamble.tex
% and nothing else. Every number in the drawing is written by
% figs/splitref_numbers.py into figs/fig-splitref.json; run that script if the
% source records ever change, and paste its "TikZ dot values" block below.
%
% STALE GENERATOR -- READ BEFORE RE-RUNNING. splitref_numbers.py:57 hardcodes
% REPRO_LINE = 0.109 and its own comment (:54-56) says 0.109 was chosen because
% it "reproduces all three printed values". That is backwards: the line was
% fitted to the prose rather than taken from the definition, and the prose value
% it was fitted to is the one that is wrong. Panel B below is now drawn at the
% resolved +0.1086 line the document actually defines (see THE LINE). Re-running
% the script as it stands will silently revert this file. Set REPRO_LINE
% = 0.1086 there first, and regenerate figs/fig-splitref.json, which still
% records the superseded 0.109 reconstruction.
%
% v1's version sat inline in Sections/05-Limitations.tex at lines 193-249.
% Four things were wrong with it and are fixed here.
%
%   1. IT PAID FOR FULLWIDTH AND DID NOT USE IT. The old float opened
%      \begin{fullwidth} -- 174mm, and by the preamble's own rule the page
%      carrying one may then carry no margin note -- and drew 150.2mm. It shows
%      six numbers, all six of which its own caption already states, so it was
%      buying 32mm of extra measure to repeat the caption. This draws inside the
%      142mm body measure (see MEASURED below) and that page keeps its margin
%      column.
%
%   2. THE RIGHT-HAND BRACE WAS A DEGENERATE GLYPH. Three values spanning
%      0.956-0.961 on an axis running 0.82-0.98 collapse a
%      decorations.pathreplacing brace to about 1.5pt of path, shorter than its
%      own 3.5pt amplitude, and pgf draws the two halves overlapping: it printed
%      as a stray chevron beside "0.006". A brace that renders as a chevron is
%      worse than no brace. Both spans are now square brackets -- a form that
%      degrades gracefully, so the SAME mark is used in both panels and the
%      reader compares 47.8mm of bracket against 1.8mm of bracket directly.
%      Correspondingly the three converging leaders on the right panel's three
%      near-coincident dots are gone: a single block set hard against the
%      cluster lists all three values in the dots' own top-to-bottom order, so
%      it reads as the cluster magnified rather than as a legend.
%
%   3. NEITHER PANEL SAID WHAT WAS ON THE VERTICAL AXIS. The caption called it
%      "the reference" and the drawing named nothing. There is now a rotated
%      axis title naming the estimator and its frame.
%
%   4. THE TRUNCATION WAS SILENT. The axis runs 0.82-0.98 -- 16 points of a
%      100-point scale -- and that truncation is most of why the 0.132 spread
%      looks dramatic. Both panels share the scale, so the before/after
%      comparison is honest, but the reader has to be told. The axis is drawn
%      BROKEN, with 0.50 (chance, the bottom of the NIR scale's meaningful
%      range) marked below the break, and the caption states the range in words.
%
% GEOMETRY. x and y are both 1mm, so every coordinate below is a millimetre and
% the numbers can be checked against the MEASURED reading. \yv maps a reference
% value onto the axis: \lo..\hi spans \H mm.
%
%   spine            x = 7.5, y = 0 (0.82) .. 58 (0.98)
%   break            gap y = -4 .. 0, two slashes across it, stub -8 .. -4
%   panel A          bracket x = 13, dots x = 36, names from 38.5, values to 66
%   divider          x = 68.5
%   panel B          bracket x = 88, dots x = 91.5, value block from x = 94.5
%   gridlines        stop at x = 94, clear of the value block, uniform length
%
% The two brackets sit on OPPOSITE sides of their dot columns and that is
% deliberate. Panel A's is 47.8mm tall and reads as a measure from anywhere;
% panel B's is 1.8mm and would read as a stray mark if it were the 23mm from
% its dots that panel A's is, so it is set hard against them and its label is
% thrown outward instead. Same mark, same weight, same colour; only the side
% differs, and the two lengths stay directly comparable across the divider.
%
% The two lower dots of panel A sit 0.0127 apart, which is 4.62mm here, against
% a \scriptsize line box of 3.35mm -- so both labels sit at their own dot's
% height and NO leader is needed on the left panel either. That clearance is
% why \H is 58 and not the 52 the old version used.
%
% MEASURED, not estimated. figs/_harness_splitref.tex \input s this file against
% the real preamble, captures the float into a box with \caption and \label
% gobbled, and \typeout s \wd. Reading:
%   tikzpicture  377.33pt = 132.90mm wide, 220.73pt = 77.74mm tall
%   \textwidth   404.03pt = 142.00mm
%   fills 93.4% of the measure, 26.70pt = 9.39mm of slack.
% The build's overfull-hbox gate is 0; this figure contributes none. The
% fullwidth version it replaces drew 150.2mm against a 174mm allowance.
%
% NUMBERS. Panel A is read straight from CANONICAL_NUMBERS.json -> ceilings{}.
% Panel B is reconstructed from re_v3.json -> records[] at repro_cos > 0.1086;
% no file in RESULTS_cluster holds those three values. figs/fig-splitref.json
% carries the full provenance, and is STALE: it still describes the superseded
% 0.109 reconstruction (299 / 394 / 111 kept, train 0.955).
%
% THE LINE. 04c-reencoding.tex:212-213 defines it as the measured null's 95th
% percentile, "+0.109, resolving to +0.1086", and the document is consistent at
% 0.1086, not at 0.109:
%   - 04c:220-221  "Of the 300 scored training conditions ... none falls at or
%                  below the resolved +0.1086." The train minimum repro_cos in
%                  re_v3.json is 0.10892294 -- above 0.1086, BELOW 0.109. At
%                  > 0.109 the figure dropped exactly the condition the text
%                  says is not dropped.
%   - 04c:959      "499 conditions sit below the +0.1086 line" on unseen_combo,
%                  i.e. 396 of 895 kept. The source count of records under
%                  0.1086 is exactly 499; under 0.109 it is 501, giving 394.
%   - 04c:957-958  "train is by construction the reliability-surviving
%                  population", so imposing the line on train must not move
%                  train. At 0.1086 it does not: panel B's train is the same
%                  0.9555006 as panel A's, to the last digit. At 0.109 it moved.
% Recomputed from re_v3.json at repro_cos > 0.1086: kept 300 / 396 / 111, means
% 0.9555006 (train), 0.9605642 (unseen_combo), 0.9597734 (unseen_drug).
%
% TWO NUMBERS THIS FIGURE PRINTS THAT THE PROSE CONTRADICTS. Neither is
% arbitrated here; both are reported upward. The figure does not invent a third
% reading of either.
%   (a) train, panel B. Drawn and printed 0.956, the value the resolved line
%       gives. The prose still prints 0.955 at 04c-reencoding.tex:883 and
%       05-Limitations.tex:224, both inherited from the 0.109 reconstruction.
%       Those two sites need the same correction; they are outside this file.
%   (b) the right-hand spread, printed 0.006. Under the old 0.109 line that was
%       max(rounded) - min(rounded) = 0.961 - 0.955, the raw spread being
%       0.00517. Under the resolved line BOTH conventions give 0.005: the raw
%       spread is 0.0050636, and max(rounded) - min(rounded) = 0.961 - 0.956.
%       So 0.006 is no longer derivable at the line the document defines. It is
%       nevertheless what the prose prints in three places
%       (04c-reencoding.tex:884, 05-Limitations.tex:225 and 03-Apparatus.tex:378,
%       the last a table note in the apparatus chapter), so it is left standing
%       here rather than silently changed to 0.005 -- changing it would put the
%       figure at odds with the prose in three places, and the spread is a
%       claim. FLAGGED, NOT FIXED.
% The left-hand 0.132 is the raw spread rounded and is unaffected.
% ===========================================================================
\begin{figure}[htbp]
\centering
\begin{tikzpicture}[x=1mm, y=1mm, font=\small]
  % --- the constants every coordinate is derived from ----------------------
  \def\lo{0.82}\def\hi{0.98}\def\H{58}
  \newcommand{\yv}[1]{{(#1-\lo)/(\hi-\lo)*\H}}
  \def\ax{7.5}          % spine
  \def\bra{13}\def\dota{36}\def\laba{38.5}\def\vala{66}
  \def\divx{68.5}
  \def\brb{88}\def\dotb{91.5}\def\blk{94.5}
  \def\srf{2}           % bracket serif length, both panels

  \tikzset{
    grid/.style   ={draw=rulegrey!45, line width=0.3pt},
    spine/.style  ={draw=rulegrey, line width=0.5pt},
    tick/.style   ={font=\scriptsize, text=quietink},
    dot/.style    ={fill=ink},
    % text height/depth normalised so a \texttt name and a maths value set at
    % the same y sit on the same baseline; without it Inconsolata's box and the
    % maths box differ by 0.7pt and the three rows print visibly askew.
    lbl/.style    ={font=\scriptsize, text=ink,
                    text height=1.55ex, text depth=0.30ex},
    blk/.style    ={font=\scriptsize, text=ink},
    span/.style   ={draw=accent, line width=0.5pt},
    strong/.style ={font=\scriptsize\bfseries, text=accent}}

  % ---- 1. the shared scale ------------------------------------------------
  % One axis serves both panels; that is the argument, so it is drawn once.
  \foreach \v in {0.82,0.86,0.90,0.94,0.98}{
    \draw[grid] (\ax,\yv{\v}) -- (94,\yv{\v});
    \node[tick, anchor=east] at (\ax-1.5,\yv{\v}) {$\v$};}

  \draw[spine] (\ax,0) -- (\ax,\H);
  % the break, and chance below it. 0.50 is the foot of the NIR scale's
  % meaningful range and it is 32 points below anything drawn, so it is marked
  % past a break rather than pretended onto the plotted span.
  \draw[spine] (\ax,-8) -- (\ax,-4);
  \draw[spine] (\ax-1.3,-3.0) -- (\ax+1.3,-2.2);
  \draw[spine] (\ax-1.3,-1.8) -- (\ax+1.3,-1.0);
  \node[tick, anchor=east] at (\ax-1.5,-8) {$0.50$};
  \node[tick, anchor=west] at (\ax+1.5,-8) {chance};

  \node[rotate=90, anchor=south, font=\scriptsize, text=quietink, align=center]
    at (-1,29) {within-well split-half reference\\(residual frame, $\NIR$)};

  \draw[draw=rulegrey!60, line width=0.4pt] (\divx,-1) -- (\divx,60);

  \node[strong, anchor=south] at (39.5,61.5) {as the splits are actually scored};
  \node[strong, anchor=south] at (94,61.5)
    {with the training line imposed on all three};

  % ---- 2. panel A: the reference as each split is actually scored ---------
  % values: CANONICAL_NUMBERS.json -> ceilings{}
  % Name anchored west at \laba, value anchored EAST at \vala, so the three
  % values line up in a column despite the three split names differing by six
  % characters. Left-setting the whole label leaves them ragged.
  \fill[dot] (\dota,\yv{0.955501}) circle (1.4pt);
  \fill[dot] (\dota,\yv{0.836477}) circle (1.4pt);
  \fill[dot] (\dota,\yv{0.823720}) circle (1.4pt);
  \node[lbl, anchor=west] at (\laba,\yv{0.955501}) {\texttt{train}};
  \node[lbl, anchor=west] at (\laba,\yv{0.836477}) {\texttt{unseen\_drug}};
  \node[lbl, anchor=west] at (\laba,\yv{0.823720}) {\texttt{unseen\_combo}};
  \node[lbl, anchor=east] at (\vala,\yv{0.955501}) {$0.956$};
  \node[lbl, anchor=east] at (\vala,\yv{0.836477}) {$0.836$};
  \node[lbl, anchor=east] at (\vala,\yv{0.823720}) {$0.824$};

  \draw[span] (\bra+\srf,\yv{0.955501}) -- (\bra,\yv{0.955501})
           -- (\bra,\yv{0.823720}) -- (\bra+\srf,\yv{0.823720});
  % 0.889610 and 0.958032 below are the MIDPOINTS of the two spans, used to
  % place each label at the middle of its own bracket. They are geometry, not
  % data, and appear nowhere in fig-splitref.json.
  \node[strong, anchor=west] at (\bra+\srf+1.5,\yv{0.889610}) {spread $0.132$};

  % ---- 3. panel B: the same reference, training line imposed on all three --
  % values: reconstructed from re_v3.json -> records[] at repro_cos > 0.1086,
  % the resolved line of 04c-reencoding.tex:212-213. Kept 300 / 396 / 111.
  % train is 0.955501 -- byte-identical to panel A's train above, because all
  % 300 train conditions clear the line. That identity is the point: the line
  % selects train, so imposing it on train cannot move train.
  \fill[dot] (\dotb,\yv{0.960564}) circle (1.4pt);
  \fill[dot] (\dotb,\yv{0.959773}) circle (1.4pt);
  \fill[dot] (\dotb,\yv{0.955501}) circle (1.4pt);

  % Bracket hard against the dots, opening toward them; label thrown left.
  \draw[span] (\brb+\srf,\yv{0.960564}) -- (\brb,\yv{0.960564})
           -- (\brb,\yv{0.955501}) -- (\brb+\srf,\yv{0.955501});
  % 0.006 is the prose's number, not this panel's: see (b) in the header. The
  % bracket it labels spans 0.0050636, i.e. 0.005.
  \node[strong, anchor=east] at (\brb-1.5,\yv{0.958032}) {spread $0.006$};

  % Three dots 1.8mm apart cannot carry three leaders, and a leader short enough
  % to reach a block set beside them is indistinguishable from a hyphen. So: no
  % leader. One block, hard against the cluster with nothing else within 25mm of
  % it, listing all three in the dots' own top-to-bottom order, so it reads as
  % the cluster magnified rather than as a legend.
  \node[blk, anchor=west] at (\blk,\yv{0.958032})
    {\begin{tabular}{@{}l@{\hspace{2mm}}l@{}}
       \texttt{unseen\_combo} & $0.961$\\
       \texttt{unseen\_drug}  & $0.960$\\
       \texttt{train}         & $0.956$
     \end{tabular}};
\end{tikzpicture}
\caption[The split references, before and after matching]{\textbf{The three
splits do not differ in recoverable signal; they differ in which conditions were
allowed into them.} The vertical axis is the within-well split-half precision
reference --- one half of a condition's cells scored against the other half's
truth, in the residual frame, on the $\NIR$ scale --- and it is
\emph{truncated} to $0.82$--$0.98$, a sixth of the scale; chance is $0.50$,
marked below the axis break and far under everything drawn. Both panels share
that truncated scale, so the two spans may be read against each other. Left,
the reference as each split is actually scored, $0.956$ on \texttt{train}
against $0.824$ and $0.836$ on the held-out splits, a spread of $0.132$. Right,
the same reference with the resolved $+0.1086$ reliability line imposed on all
three, $0.956$, $0.961$ and $0.960$, a spread of $0.006$; at this scale the
three sit inside a bracket two millimetres tall, so they are listed beside the
cluster in
their own order rather than labelled in place. The left-hand spread is a
selection artefact. The markers are bare: no interval is carried for this
reference in the source record, and none is drawn rather than invented.}
\label{fig:splitref}
\end{figure}


\emph{Those conditions carry less signal} is therefore not a reason to discount
a held-out shortfall, and this thesis nowhere uses it as one. What survives is
narrower and binding. The matched-population rescore covered the precision
reference only. The model's own arms were never rescored, so what the held-out
shortfall becomes once the populations agree is not reported here, and the
training-exposure contrast of \cref{sec:res-generalisation} inherits the
mismatch.

% PLACEMENT ONLY. Eight lines and a framed header, and it was splitting five and
% three across the page turn, so the header sat with half its own text and the
% rest arrived under the running head as an unlabelled fragment. 11 lines is the
% box entire.
\needspace{11\baselineskip}%
\begin{scopelimit}[discharges the limit carried from \cref{sec:res-lookup}]
Three selection rules are in force at once, and no result computed under one of
them may be set beside a result computed under another. A raw score read across
the three splits is not answering the same question on each, because the splits
were filtered by different rules. An absolute $\NIR$ computed with the
comparison set grouped by cell line is not comparable with one computed within
plate. And a per-drug lookup is the bar for a drug that was \emph{seen} in
training and is being asked about a new context: no lookup exists for a drug
that never appeared in training, so that regime keeps its own bar, set in
\cref{sec:res-scope}.
\end{scopelimit}

% ===========================================================================
\section{What the unseen-drug arm licenses}

Two defects in the unseen-drug arm must be stated. A third, the comparator, was
found and corrected, and \cref{sec:res-generalisation} reports the arm against
the neutral stratum throughout.

\paragraph{The design leaks the drug's class} Holding out a drug does not
remove everything the model knows about it. The prompt contains
\texttt{Mechanism: \{moa\}}, so a held-out drug still arrives with its class
label. The correct design is leave-one-\emph{mechanism}-out.

The leak is measured, and the measurement is not clean. Scrambling the
\texttt{Mechanism} field alone moves the model's $\NIR$ by
\ci{+0.0303}{-0.0078}{+0.0685} ($n = \num{738}$, clustered on drug and well), an
interval containing zero. That understates the channel; it does not bound it,
because this row's swap partner scores $0.520$, above chance, and a comparator
above chance suppresses the gap (\cref{tab:fieldattribution}, Panel~B). A
mechanism-only lookup is the same measurement from the other side and equally
inconclusive. Pooled over the $n = \num{218}$ conditions on which both are
scored, model $-$ \texttt{moa\_lookup} $=$ \ci{+0.044}{-0.048}{+0.135},
clustered two-way on cell line and well, spans zero, so the comparison settles
nothing in either direction.

On \texttt{unseen\_drug}, the split where it would matter, the mechanism lookup
is ahead of the model by $0.099$ on $44$ conditions; in the contrast's reported
direction, model minus lookup is \ci{-0.099}{-0.222}{+0.008}, clustered on cell
line, spanning zero and far too wide to settle anything. That measurement has no
power at exactly the point where the design defect bites. It is not evidence
that the leak is doing work. The design is wrong, the arm cannot say what the
wrongness cost, and its proper use is as an instrument control.

\paragraph{It is underpowered by construction} The half-width on the
\texttt{unseen\_drug} gap is $\pm 0.052$ clustered two-way on cell line and
well, and $\pm 0.073$ clustered on the drug; the smallest effect detectable at
$80\%$ power is $\approx 0.079$. The largest effect the mechanism channel could
deliver is the distance from chance to \texttt{moa\_lookup}, which scores
$0.552$ pooled across splits against a chance value of $0.500$, so $\approx
0.052$. The arm \emph{cannot} detect the largest effect it could plausibly be
measuring (\cref{fig:power}).

% ===========================================================================
% power.tex -- the unseen-drug arm's detection floor against the largest
% effect the mechanism channel could plausibly deliver.
% \input by Sections/05-Limitations.tex, in "It is underpowered by
% construction". Compiles against thesis_v2/preamble.tex and nothing else;
% no \includegraphics, no external file read at build time.
%
% WHY THIS FIGURE EXISTS. The passage it sits in makes its argument out of
% four numbers delivered in four separate sentences (0.052, 0.073, 0.079,
% 0.052-as-largest-effect). The conclusion -- the arm cannot detect the
% largest effect it could plausibly be measuring -- only follows once all
% four are on one scale. On one scale it is a picture: the bar ends before
% the threshold begins.
%
% ---------------------------------------------------------------------------
% EVERY NUMBER DRAWN, ITS SOURCE, AND THE THESIS SENTENCE IT MUST MATCH.
% Sibling power.json records the same list machine-readably.
%
%  threshold (dashed accent rule)      0.07897560995161473
%      RESULTS_cluster/CANONICAL_NUMBERS.json
%          .detection_limit.detectable_at_80pc_power
%      (same object carries .split = "unseen_drug", .n = 199,
%       .observed = 0.45889838436740604, .half_width = 0.05528292696613031;
%       the observed NIR and that half_width are NOT drawn -- this is a
%       difference axis and 0.4589 is an absolute NIR.)
%      Sections/05-Limitations.tex:310-311 "the smallest effect detectable
%      at $80\%$ power is $\approx 0.079$".                        MATCHES.
%
%  bar length                          0.05201514916327314
%      = 0.5520151491632731 - 0.500
%        RESULTS_cluster/CANONICAL_NUMBERS.json .means_all.moa_lookup
%        = RESULTS_cluster/re_v3.json .means.moa_lookup   (identical to 16 sf)
%        chance 0.500 is a definition: both frames are two-alternative
%        retrievals (preamble/frames.tex, Sections/03-Apparatus.tex).
%      Sections/05-Limitations.tex:311-314 "The largest effect the mechanism
%      channel could deliver is the distance from chance to
%      \texttt{moa\_lookup}, which scores $0.552$ pooled across splits
%      against a chance value of $0.500$, so $\approx 0.052$."      MATCHES.
%      0.552 and 0.500 are absolute NIR values and are therefore stated in
%      the caption only; the drawing carries the difference and nothing else.
%
%  point estimate (filled marker)     -0.03148175705606285
%      CANONICAL_NUMBERS.json .splits.unseen_drug.gap
%      = re_v3.json .means.generalization.unseen_drug.gap
%      Sections/05-Limitations.tex:322 \ci{-0.0315}{-0.1049}{+0.0419};
%      Sections/04c-reencoding.tex:791 table row, same estimate twice. MATCHES.
%
%  inner span (thick)  -0.08358077892636565 .. +0.020617264814239943
%      CANONICAL_NUMBERS.json .splits.unseen_drug.gap_ci_line
%      half-width = (0.020617264814239943 + 0.08358077892636565)/2
%                 = 0.052099021870302796  -> 0.052
%      Sections/05-Limitations.tex:308-310 "$\pm 0.052$ clustered two-way on
%      cell line and well". The key "gap_ci_line" is the JSON's own shorthand;
%      tab:generalisation's note says the interval is two-way cluster-robust
%      over cell lines and wells on min(G_line,G_well)-1 = 27 df, so the
%      figure uses the prose's two-way wording.                     MATCHES.
%
%  outer span (thin)   -0.10487258422550776 .. +0.04190907011338206
%      CANONICAL_NUMBERS.json .splits.unseen_drug.gap_ci_drug
%      half-width = (0.04190907011338206 + 0.10487258422550776)/2
%                 = 0.07339082716944491   -> 0.073
%      Sections/05-Limitations.tex:309-310 "$\pm 0.073$ clustered on the
%      drug"; the same interval is printed at 05:322 and at
%      04c-reencoding.tex:791.                                      MATCHES.
%
%  n = 199
%      CANONICAL_NUMBERS.json .splits.unseen_drug.n
%      Sections/04c-reencoding.tex:726 "$\num{199}$ unseen-drug" and the
%      tab:generalisation row "$199$ ($40$; $10$; $28$)".            MATCHES.
%
%  Both intervals are exactly symmetric about the point estimate: the
%  midpoint of gap_ci_line and the midpoint of gap_ci_drug each equal
%  -0.031481757056062853, i.e. .splits.unseen_drug.gap to 16 significant
%  figures. So "point estimate +- half-width" is the intervals themselves and
%  not a symmetrisation of them; nothing is being smoothed here.
%
% DELIBERATELY ABSENT, per the brief: any absolute NIR value inside the
% drawing (0.4589, 0.552, 0.500, 0.8365); any other split; any p-value; any
% legend (both spans are labelled where they end); any gridline.
%
% ---------------------------------------------------------------------------
% GEOMETRY. x is in NIR-difference units at x=383mm, so every coordinate in
% the source below is the datum itself and cannot drift from it by arithmetic
% done here. y is in mm.
%
%   axis          -0.12 .. +0.12   ->  -45.96mm .. +45.96mm   (91.92mm)
%   left column   44mm of flush-right label + 3mm gutter, east edge -48.96mm
%   total width   44 + 3 + 45.96 + 45.96 = 138.92mm, inside the 142.0mm
%                 measure (404.03pt, preamble.tex:793). Measured on the
%                 harness build: see _harness_power.log "HARNESS tikz width".
%   bar tip       0.05201515 * 383 = 19.92mm
%   threshold     0.07897561 * 383 = 30.25mm
%   -> the bar stops 10.33mm short of the rule. That clearance is the whole
%      argument, so it is stated here as a measured length and not judged by
%      eye; at 142mm it is about 7 per cent of the measure and reads at once.
%
%   rows, y in mm   bar      33.5 (rectangle 30.7..36.3)
%                   estimate 13.5 (inner caps 11.1..15.9, outer 12.0..15.0)
%                   axis      0.0, ticks to -1.4, axis title at -7.6
%   detection rule  -1.4 .. 42.5, its label sitting on top of it at 42.6
%   zero rule       -1.4 .. 16.4 -- see note at the draw itself
%
%   MEASURED on the harness build (_harness_power.log):
%     HARNESS textwidth  = 404.02913pt = 142.00mm
%     HARNESS tikz width = 395.70499pt = 139.07mm   <- 2.93mm inside the
%     HARNESS tikz height= 173.09203pt =  60.83mm      measure, 0 overfull
%   and in the whole-thesis build the only Overfull \hbox in main.log is the
%   pre-existing 1.35pt one in figs/fig-frame.tex, not this file.
%
% COLOUR. accent (#0B5394) is spent on exactly two marks, the bar and the
% detection floor, because those two are the comparison the passage makes.
% The observed estimate is ink, the axis and the zero rule rulegrey, all
% descriptive text quietink and every numeral ink. Nothing is told apart by
% colour that is not also told apart by position, weight or shape: the two
% nested spans differ in thickness and in length, not in hue, and the figure
% survives greyscale printing.
% ===========================================================================
\begin{figure}[htbp]
\centering
\begin{tikzpicture}[
    x=383mm, y=1mm, font=\footnotesize,
    rowlbl/.style  ={text=quietink, align=flush right, text width=44mm,
                     inner sep=0pt, anchor=east, xshift=-3mm},
    ticklbl/.style ={font=\scriptsize, text=quietink, inner sep=0pt,
                     anchor=north},
    axtitle/.style ={font=\scriptsize, text=quietink, inner sep=0pt,
                     anchor=north},
    note/.style    ={text=quietink, inner sep=0pt}]

  % ---- 1. the difference axis ---------------------------------------------
  % Zero, not 0.50, is the reference: this axis carries differences, so the
  % chance level lives inside each quantity's definition and not on the rule.
  \draw[rulegrey, line width=0.5pt] (-0.12,0) -- (0.12,0);
  \foreach \t/\lab in {%
      -0.10/$-0.10$, -0.05/$-0.05$, 0/$0$, 0.05/$+0.05$, 0.10/$+0.10$}{%
    \draw[rulegrey, line width=0.5pt] (\t,0) -- (\t,-1.4);
    \node[ticklbl] at (\t,-2.1) {\lab};}
  \node[axtitle] at (0,-7.6) {difference in $\NIR$\enspace($0$ = no difference)};

  % ---- 2. zero -------------------------------------------------------------
  % The rule stops at 16.4, just under the inner span's label. It does not
  % need to reach the bar row: the bar's own left edge IS x=0 and carries the
  % position up on its own. Running it higher would have to be struck through
  % the "+-0.052" label, which is 46mm wide and must end flush with the inner
  % span's right cap at +0.0206, so it cannot avoid crossing zero. A knocked-
  % out rule reads as a broken line; a shorter one reads as a rule.
  \draw[rulegrey, line width=0.4pt] (0,-1.4) -- (0,16.4);

  % ---- 3. the detection floor ---------------------------------------------
  \draw[accent, line width=0.9pt, dash pattern=on 2.2pt off 1.6pt]
    (0.07897561,-1.4) -- (0.07897561,42.5);
  \node[note, anchor=south east, align=flush right, xshift=-1.0mm]
    at (0.07897561,42.6)
    {smallest effect this arm can\\detect at $80\%$ power,
     \textcolor{ink}{$0.079$}};

  % ---- 4. the largest effect the mechanism channel could deliver ----------
  % It starts at zero because it IS a distance from a null, and it ends
  % where it ends. The dotted arrow is the shortfall: it is drawn without a
  % number on it because the thesis states 0.052 and 0.079 and does not
  % state their difference.
  \fill[accent!20]                (0,30.7) rectangle (0.05201515,36.3);
  \draw[accent, line width=0.7pt] (0,30.7) rectangle (0.05201515,36.3);
  \draw[quietink, line width=0.6pt, dotted, line cap=round,
        -{Latex[length=1.5mm,width=1.3mm]}]
    ([xshift=1.4mm]0.05201515,33.5) -- ([xshift=-0.8mm]0.07897561,33.5);
  \node[rowlbl] at (-0.12,33.5)
    {largest effect the\\mechanism channel\\could deliver\\[0.35ex]
     \textcolor{ink}{$0.052$}};

  % ---- 5. the arm's own estimate, with both cluster keys ------------------
  % Two nested spans, same point estimate, different unit of generalisation.
  % Thin/long = clustered on drug; thick/short = clustered two-way on cell
  % line and well. Told apart by weight and by extent, never by colour.
  % Each label ends FLUSH with the span end it names -- the inner one at its
  % right cap, above; the outer one at its left cap, below -- so the two are
  % separated diagonally and neither needs a legend or a leader line.
  \draw[ink, line width=0.5pt] (-0.10487258,13.5) -- (0.04190907,13.5);
  \draw[ink, line width=0.5pt] (-0.10487258,12.0) -- (-0.10487258,15.0);
  \draw[ink, line width=0.5pt] ( 0.04190907,12.0) -- ( 0.04190907,15.0);
  \draw[ink, line width=1.3pt] (-0.08358078,13.5) -- (0.02061726,13.5);
  \draw[ink, line width=1.0pt] (-0.08358078,11.1) -- (-0.08358078,15.9);
  \draw[ink, line width=1.0pt] ( 0.02061726,11.1) -- ( 0.02061726,15.9);
  \fill[white] (-0.03148176,13.5) circle (2.1pt);
  \fill[ink]   (-0.03148176,13.5) circle (1.6pt);

  \node[note, anchor=south east, align=flush right] at (0.02061726,16.6)
    {$\pm 0.052$ clustered two-way\\on cell line and well};
  % broken over two lines, not one, so that its right edge stops at -0.026
  % and the label clears the zero rule instead of being knocked out of it.
  \node[note, anchor=north west, align=flush left] at (-0.10487258,10.4)
    {$\pm 0.073$ clustered\\on drug};

  \node[rowlbl] at (-0.12,13.5)
    {observed \texttt{unseen\_drug} gap\\($n = 199$)\\[0.35ex]
     \textcolor{ink}{$-0.0315$}};
\end{tikzpicture}
\caption[The detection floor sits above the largest plausible
effect]{\textbf{The bar ends before the threshold begins.} Every quantity
here is a difference in $\NIR$ on one scale; no absolute $\NIR$ is plotted.
The dashed rule is the smallest effect the \texttt{unseen\_drug} arm can
detect at $80\%$ power, $0.079$, on the $n = \num{199}$ conditions that arm
scores. The bar is the largest effect the mechanism channel could deliver ---
the distance from chance to the mechanism-only lookup, which scores $0.552$
pooled across splits against a chance value of $0.500$, so $\approx 0.052$ ---
and it stops short of the rule. Beneath it is the arm's own estimate: model
$\NIR$ minus the geometry-defined \texttt{orth} scramble comparator on the
same conditions, $-0.0315$, drawn with the thick span the $\pm 0.052$
half-width clustered two-way on cell line and well and the thin span the
$\pm 0.073$ half-width clustered on the drug. Both cover zero, and the
half-width on the tighter of the two is itself the size of the bar. The arm
is licensed as an instrument control and is not evidence that unseen-drug
transfer is zero: an effect of the largest size the mechanism channel could
deliver would sit inside that same interval, so what it supplies is the
absence of a contradiction rather than a confirmed null.}
\label{fig:power}
\end{figure}


That cuts both ways, but asymmetrically. As a \emph{control} the arm is not
falsified, since the estimate \ci{-0.0315}{-0.1049}{+0.0419}, clustered on drug,
covers zero. It does not sit \emph{at} zero, being negative and of the same
order as the largest effect the arm could plausibly be measuring, so what it
supplies is the absence of a contradiction. That is why the arm validates the
comparator in \cref{sec:res-generalisation} but cannot argue that unseen-drug
transfer is zero. ``Unseen'' is narrow in a second way too. Pythia-1b's
pretraining is natural-language text in which drug names occur, and the
prompt still supplies \texttt{Mechanism}.

\begin{scopelimit}[carried into \cref{sec:conclusions}]
The \texttt{unseen\_drug} arm is licensed as an instrument control. It is not
evidence that unseen-drug transfer is zero. Its estimate covers zero, but an
effect of the largest size the mechanism channel could deliver would sit inside
that same interval, so what it supplies is the absence of a contradiction rather
than a confirmed null. ``Unseen'' means held out from Tahoe fine-tuning and
nothing wider than that.
\end{scopelimit}

% ===========================================================================
\section{What would settle it}

A limitation earns its place when something could be run to remove it.

\paragraph{The arm to build carefully} The obvious continuation of
\cref{sec:res-scope} is a channel-conditioned prompt that appends protein
targets to the instruction, retrains, and evaluates on a leave-one-mechanism-out
split. We do not recommend building it eagerly.

\cref{sec:res-field} does not resolve \emph{which} span of the prompt is read.
Swapping the drug \emph{name} alone moves the prediction by
\ci{+0.1026}{+0.0625}{+0.1427} ($n = \num{1394}$); swapping the \emph{mechanism}
field alone moves it by \ci{+0.0303}{-0.0078}{+0.0685} ($n = \num{738}$), both
clustered on drug and well. Both swap to an opposite-signature partner, which is
not a null, so each gap carries a comparator term running in opposite directions
across the rows, inflating the name arm, whose partner scores $0.434$, and
suppressing the mechanism arm, whose partner scores $0.520$. Read as each row's
distance from chance the ordering inverts ($+0.0367$ on the name row against
$+0.0506$ on the mechanism subset), so the split between the two spans is
unresolved.

Pending a re-run against the neutral partner, the case against a
target-conditioned prompt rests not on demonstrated mechanism-blindness but on
the size of the readout deficit. On the common support the model covers $13.9\%$
of the distance from chance to the within-well bar that a context-blind lookup
covers entirely, and adding a field to a prompt does not touch that deficit.

Worth building instead is an intervention that makes the drug \emph{causal}
rather than merely present, a learned per-drug modulation of the transformation
from control to treated, so the drug parameterises the computation instead of
sitting in the residual stream as an input the readout can round away. That is a
research architecture, outside the scope of what remained here.

\paragraph{Two channel questions that remain open} Target and mechanism overlap,
so whether target adds anything \emph{over} mechanism is unresolved, and it
matters: the prompt already exposes mechanism, and whether the model reads that
field is itself unsettled. Second, the gate of \cref{sec:res-scope}
measures \emph{similarity transfer}, an unseen drug resembling a seen drug with
the same target. It does not test a \emph{compositional rule}, that a named
target's own pathway moves, which would apply to any drug whose target is named,
including one with no similar drug in training. That claim is the stronger one,
it does not reduce to a lookup, and it is untested.

Both sat behind a defect in the gate itself. The partner pool was not restricted
to training drugs, so a held-out drug's response could be predicted from a
partner that was also held out; the gate has since been re-run with the
restriction in place.\corrmark{the gate's partner pool, restricted and re-run in
\cref{sec:res-scope}} The repair narrowed what may be claimed, and the
clustering axis narrows it again. The target, mechanism and chemistry arms
differ in support, drug set, partner count and annotation coverage, so their
estimates are descriptive and do not identify a biological ranking.

Protein target clears the relevance margin on the cell-line axis, misses it on
the drug axis against both estimand-correct nulls, clearing it there only
against the weaker count-matched one, and rests on $18$ annotated drugs of which
two are held out; mechanism misses its lower bound by $0.0021$ on the cell-line
axis and spans zero on the drug axis; chemistry's interval straddles the margin
on both. Chemistry attenuates under the different-plate construction, but that
construction also changes support and partner availability, so the attenuation
cannot be assigned causally to co-plating. On its restricted support, same-target
retrieval is the strongest observed lead, and stands only as a direction worth
testing. The binding limitation is coverage: $18$ drugs cannot settle it, and
the experiment that would is a target-annotated panel broad enough to cluster
on.

\paragraph{Unseen cell lines} A cell line never seen with any drug is a
different generalisation axis from the held-out wells tested here, and one the
plating design does not compromise in the same way. It asks whether the control
cell's state is used \emph{compositionally} with the
drug.\seealso{\cref{sec:res-transfer}: $\Tcoef$, and what its complement does
not measure}

The transfer coefficient scopes it only loosely. \ci{44.7\%}{40.7\%}{48.6\%} of
the similarity scale $\Tcoef$ is measured on is not shared across the compared
contexts, and no structure was recovered in that unshared part by tests that do
not settle it. That figure is not a drug$\times$cell-line interaction share,
since it mixes cell line, dose, well and estimator (\cref{sec:res-transfer}), so
it bounds what a per-drug constant leaves on the table without establishing that
a new cell line is the axis the remainder lies along.

No consistent per-line direction was found, so the achievable target on a new
cell line plausibly reduces to recalling the drug's per-drug constant, which a
training-only lookup already does at $0.890$ on held-out wells, the nearest
measured analogue of a new context ($0.913$ pooled, inflated on trained
conditions by the well advantage). The axis is less promising than it appeared,
though the argument is weaker than a coefficient-based one.

\paragraph{Objectives that were considered and rejected} Three training
objectives were designed and then set aside on their own mechanics, recorded so
they are not re-litigated.
\begin{itemize}
  \item \textbf{Reinforcement learning with a contrastive reward.} Its
        per-sequence reward keeps the token-dilution ratio out of the gradient,
        attacking a bottleneck the residual encoding has already lifted ($34.6
        \rightarrow 118.2$ differing tokens at cell-line scope, $120.3$ at the
        plate scope the trained build uses). Its reward optimises similarity to
        the true residual, which \texttt{drug\_lookup} already achieves at
        $0.913$ pooled, $0.890$ on held-out wells. Its best case is convergence
        to the lookup.
  \item \textbf{A differentiable profile-level contrastive loss.} The natural
        construction, $\textrm{score}(g) = \sum_i P(\textrm{token}_i =
        g)\,w(i)$, presumes gene $\approx$ token, but $86.9\%$ of the $\num{946}$
        panel genes share a first sub-token under this tokeniser (mean $3.254$
        BPE tokens per gene, only $8$ single-token). Under teacher forcing the
        target \emph{is} the drug-specific residual, so the true prefix
        identifies the drug within a few symbols and the gradient vanishes at
        the readout solution the loss exists to improve on.
  \item \textbf{Loss re-weighting toward differentially expressed genes.}
        A re-weighted cross-entropy against a single target contains no
        negatives, so nothing forces one drug's prediction to differ from
        another's. It leaves the discrimination failure where it was, and that
        failure is
        what the measurements establish, whether or not the readout is also
        blind to direction specifically, which \cref{sec:res-mechanism} declines
        to establish.
\end{itemize}

\paragraph{One model size, and the smallest one} Every result here comes from a
1B-parameter backbone. C2S-Scale releases checkpoints up to Gemma-2 at 27B, and
we tested none of them, so the most predictable question an examiner can ask,
whether it works if the model is bigger, has no answer in this thesis.

The constraint this work points to is not capacity. The target tokens are
overwhelmingly shared across drugs --- only $17\%$ of them differ between two
drugs in the same cell line under the standard target, against $59\%$ under the
residual encoding at cell-line scope and $60\%$ at the plate scope the trained
build uses (\cref{sec:res-encoding}) --- so very little of the training signal
can carry drug identity. That step is a hypothesis about the optimisation; no
token-level loss or gradient attribution was computed, so no fraction of the
gradient is claimed. A larger model receives the same diluted signal from the
same objective, and re-encoding lifted the effect while leaving
capacity untouched, which points away from capacity without establishing what it
points towards. But the re-encoding changed five things at once, aggregation,
control referencing, sign coding and length among them, and a budget-matched
optimal-transport target lifts the effect too, so \cref{sec:res-encoding}
declines to identify the tokenisation as the constraint that binds.

The genuinely different argument for scale is \emph{prior knowledge}. A 27B
model plausibly knows, as a fact about the world, that a named drug inhibits a
named kinase, which a 1B model may lack. \cref{sec:res-field} leaves
unresolved whether the model uses the \emph{mechanism} field it already has, at
\ci{+0.0303}{-0.0078}{+0.0685} ($n = \num{738}$) with a comparator above chance
suppressing the gap, and real pharmacological knowledge is exactly the case
where that might change. Testing it needs only a low-rank adaptation of a
mid-sized checkpoint. We did not run it, and it is the first thing we would run
with more time.

\paragraph{One epoch, and a thin slice of the atlas} Model size is pre-empted
above at some length. Data scale is the other half of that question and it is
not the same one. Fine-tuning saw $\num{620564}$ cells, $\num{600000}$ treated
and $\num{20564}$ control, drawn from the roughly one hundred million
Tahoe-100M holds, covering $106$ drugs of its $\approx\num{1100}$, for one
epoch at one learning rate (\cref{tab:trainconfig}). Under a hundredth of the
atlas was used, and no arm saw any cell twice.

The diluted-signal argument above is a claim about the objective and survives
this: a model trained on ten times the cells receives the same target strings,
of which only $17\%$ of positions differ between two drugs in one cell line
(\cref{sec:res-encoding}), so more data of the same kind does not make the drug
easier to see in the loss. The ten-epoch diagnostic separately rules out
undertraining in the optimisation sense, since validation loss is best at the
first epoch and worse at every epoch reported after it. Neither argument rules out
the possibility that a drug needs to be seen in more cellular contexts than
this sample provides before its response becomes learnable, and that
possibility is \notestablished either way. It is a cheaper experiment than a
larger backbone and it is not one we ran.

\paragraph{The starting checkpoint was never prepared for this task} We
fine-tune from \texttt{C2S-Scale-Pythia-1b-pt}, whose cell-sentence pretraining
corpus is atlas data alone, CELLxGENE and the Human Cell Atlas, with no
perturbation corpus of any kind, and whose released capabilities do not include
perturbation prediction.

That is why the held-out splits are as clean as they are. A drug withheld here
entered neither the fine-tuning data nor the generic that defines its target,
the split being assigned from metadata before any expression was aggregated, so
no Tahoe perturbation response for it can have been memorised. But the failure
we report is then a failure of \emph{this} starting point, and the fair
follow-up is whether a checkpoint already tuned for perturbation prediction
would behave differently.

Within the C2S-Scale line that question cannot be answered from what is public.
The perturbation-specialised model there is a Gemma-3 4B checkpoint that was not
released. Outside that line it is answerable, and we did not answer it.
Tahoe-x1~\cite{tahoex1} is a perturbation-trained foundation model on this same
atlas, public before this thesis was submitted. It is not a cell-sentence model,
so it does not isolate the representation variable, but it is the obvious
external check on whether a readout failure of this kind survives
perturbation-specific pretraining, and it was not run here. That is a gap in
this work, not a gap in the public record.

Every arm starts from the same base and differs only in the variable under test,
with one material exception, the consensus arm, stopped at roughly $60\%$ of an
epoch with its tier-1 $\NIR$ unscoreable (\cref{tab:trainconfig}). The
\emph{comparisons} between arms are therefore unaffected even where the absolute
level is not, and that arm's null is quoted with the qualification wherever it
is used.

\paragraph{Generality} Everything here is one backbone on one atlas in one
representation. Whether the readout-specificity failure belongs to cell-sentence
models specifically, or to autoregressive conditional generation over
transcriptomes generally, is open. The mechanistic protocol of
\cref{sec:res-used} to \cref{sec:res-mechanism} (purified subspace, removal
gate, positive control, matched-norm swap) transfers unchanged to any model
with an accessible residual stream.
